% id: 0502
% book: RPM 중학 3-2 · 05 원주각의 활용
% section: 유형 12 두 원에서 접선과 현이 이루는 각 (1)
% figure: crop:fig-0502.png
% answer: ②
% answer_source: 답지
% variant_level: 0 (원본 전사)
% 이 파일은 build.mjs 가 items.json 에서 만든다. 고칠 때는 items.json 을 고친다.
\smash{\raisebox{1.5mm}{\dmkichul{RPM 중학 3-2 0502번 · 83쪽}}}
\noindent\begin{minipage}[t]{0.58\linewidth}\vspace{0pt}\raggedright 오른쪽 그림에서 직선~$\pt{EF}$는 두 원~$\pt{O}$, $\pt{O'}$의 공통인 접선이고 점~$\pt{T}$는 접점이다. 점~$\pt{T}$를 지나는 두 직선이 두 원과 만나는 점을 각각 $\pt{A}$, $\pt{B}$, $\pt{C}$, $\pt{D}$라 할 때, 다음 중 옳지 \underline{않은} 것은?\end{minipage}\hfill\begin{minipage}[t]{0.4\linewidth}\vspace{0pt}\centering \setlength{\dmfigw}{86.2pt}\ifdim\dmfigw>\linewidth\setlength{\dmfigw}{\linewidth}\fi\includegraphics[width=\dmfigw]{figures/fig-0502.png}\end{minipage}\par
\begin{choicesv}{$\angle\pt{BAT}=\angle\pt{CDT}$}{$\angle\pt{ABT}=\angle\pt{CTE}$}{$\seg{AB}\parallel\seg{CD}$}{$\triangle\pt{ABT}\sim\triangle\pt{DCT}$}{$\seg{AB}:\seg{CD}=\seg{AT}:\seg{DT}$}\end{choicesv}
